\hypertarget{appendix-b60686ba06a1}{}
\section{Fourier analysis and spectral phase}\label{appendix-b60686ba06a1-fourier-analysis-and-spectral-phase}

\textbf{Curriculum depth:} Developed · \textbf{First seen:} XR01 · \textbf{Latest development:} University courses

\subsection{Definition and mechanism}\label{appendix-b60686ba06a1-definition-and-mechanism}

Fourier coordinates encode frequency magnitude and phase; phase places structure. DFT analyzes a finite periodic block, FFT computes that transform efficiently, and finite-window leakage differs from bin spacing. Interior one-sided amplitude bins need conjugate-pair normalization.

\subsection{Prerequisites}\label{appendix-b60686ba06a1-prerequisites}

\begin{itemize}
\tightlist
\item
  \hyperref[appendix-1077ea8fc737]{Continuous signals and discrete samples} --- Filtering/frequency analysis operates on a declared domain and its samples.
\end{itemize}

\subsection{Used by}\label{appendix-b60686ba06a1-used-by}

\begin{itemize}
\item
  \hyperref[appendix-6a7869dd79a6]{Sampling aliasing and reconstruction}
\item
  \href{../03-Visual-Computing-I/main.pdf\#nameddest=appendix-d5b3c692d9f0}{Convolution and filter footprints (03-Visual-Computing-I, appendix-d5b3c692d9f0)}
\item
  \hyperref[appendix-8b76b0efd27f]{PSF OTF and MTF}
\item
  \hyperref[appendix-b2c1846508bf]{Perceptual metrics and display calibration}
\item
  \hyperref[appendix-3c73f7862157]{Exceptional projections and dimension} \#\# Recall and next depth
\end{itemize}

Recover the two tested sinusoid amplitudes and explain leakage. Pending: FFT implementations, windows and finite-boundary estimates. Written coverage is not a claim of learner mastery or accepted numbered-curriculum progress.

\subsection{Source and navigation}\label{appendix-b60686ba06a1-source-and-navigation}

\hyperref[04-visual-computing-ii-ch042-fourier-coordinates-reveal-what-a-filter-removes]{Expanded teaching explanation} · \href{../01-Foundations-I/main.pdf\#nameddest=appendix-fd090d982a15-xr01---independent-seeing-and-representation-source}{Independent coverage (01-Foundations-I, appendix-fd090d982a15-xr01-\/-\/-independent-seeing-and-representation-source)} · \href{../02-Foundations-II/main.pdf\#nameddest=appendix-4ad47bdd269a}{Dependency Map (02-Foundations-II, appendix-4ad47bdd269a)}

\subsection{University course development}\label{appendix-b60686ba06a1-university-course-development}

Extend the existing frequency-domain treatment with visual interpretation and filtering consequences. A Fourier coefficient records how strongly a sinusoidal basis function occurs at a given frequency and phase. In images, the origin of a two-dimensional spectrum represents slow spatial variation; moving farther from it corresponds to finer variation. A broad bright center indicates low-frequency structure, while energy far from the center often represents edges and texture. Filtering multiplies each frequency by the filter response: a low-pass response suppresses fine detail, and a high-pass response emphasizes it. Aliasing folds unresolved high frequencies into lower ones because sampled sinusoids can share identical sample values. This explains moiré patterns and animated shimmer, and shows why a spectrum is useful for diagnosing sampling problems rather than only for theoretical signal processing.

\href{../03-Visual-Computing-I/main.pdf\#nameddest=03-visual-computing-i-ch003-sampling-continuous-signals-and-image-formation}{Sampling continuous signals and image formation (03-Visual-Computing-I, 03-visual-computing-i-ch003-sampling-continuous-signals-and-image-formation)} · \href{../03-Visual-Computing-I/main.pdf\#nameddest=03-visual-computing-i-ch004-image-sampling-filtering-and-efficient-reconstruction}{Image sampling, filtering, and efficient reconstruction (03-Visual-Computing-I, 03-visual-computing-i-ch004-image-sampling-filtering-and-efficient-reconstruction)} · \hyperref[04-visual-computing-ii-ch030-edges-contours-and-subpixel-localization]{Edges, contours, and subpixel localization}
